\documentclass[english,hebrew]{article}
\usepackage[HE8,T1]{fontenc}
\usepackage[utf8]{inputenc}
\setcounter{secnumdepth}{-2}
\setlength{\parindent}{0bp}
\usepackage{calc}
\usepackage{amsmath}
\usepackage{amsthm}
\usepackage{amssymb}

\makeatletter
\@ifundefined{date}{}{\date{}}
\usepackage{graphicx}
\usepackage{geometry}
\newcommand{\limi}{\underset{n\rightarrow\infty}{\lim}}
\newcommand{\limxz}{\underset{x\rightarrow x_{0}}{\lim}}
\newcommand{\limfxz}{\underset{x\rightarrow x_{0}}{\lim}f(x)}
\newcommand{\limfxm}{\underset{x\rightarrow x_{0}^{-}}{\lim}f(x)}
\newcommand{\limfxp}{\underset{x\rightarrow x_{0}^{+}}{\lim}f(x)}
\newcommand{\limxm}{\underset{x\rightarrow x_{0}^{-}}{\lim}}
\newcommand{\limxp}{\underset{x\rightarrow x_{0}^{+}}{\lim}}
\newcommand{\sqrm}{M_{m\times m}(\mathbb{F})}
\newcommand{\seqa}{(a_{n})_{n=1}^{\infty}}
\newcommand{\seqx}{(x_{n})_{n=1}^{\infty}}
\newcommand{\funcdr}{f:D\rightarrow\mathbb{R}}
\newcommand{\der}{\limxz\frac{f(x)-f(x_{0})}{x-x_{0}}}
\usepackage[all]{pdfprivacy}

\makeatother

\usepackage{babel}
\begin{document}
\newgeometry{top=1cm, bottom=2cm, left=1cm, right=1cm}
\title{אלגברה ליניארית {\beginL 2\endL} - הרצאה {\beginL 4\endL}}
\author{נעם לב}
\date{\L{$21.10.2026$}}

\maketitle
\noindent\fbox{\begin{minipage}[t]{1\columnwidth - 2\fboxsep - 2\fboxrule}%
\begin{center}
האוניברסיטה העברית \textbar{} אלגברה ליניארית ){\beginL 2\endL}( -
\L{$80135$} \textbar{} סמסטר{\small{} }{\beginL 2027\endL}{\small{} }א'
)תשפ\textquotedblleft ז א'( \textbar{} מרצה: פרופסור יבגני סטרחוב
\textbar{} מתרגל: מר לזר גוטרמן
\par\end{center}%
\end{minipage}}
\begin{center}
הסיכום הוא לא רשמי ולא עבר אישור של הסגל, על כן אין להשתמש בו )ובין
השאר במספור שלו ובשמות הטענות שבו( כמקור רשמי. אם ישנן טעויות, הן
שלי ולא של הסגל.
\par\end{center}

\medskip{}

\tableofcontents{}

\newpage{}

\part{ערכים עצמיים ולכסון מטריצה )המשך(}

\section{דוגמה - חישוב ערכים עצמיים, וקטורים עצמיים ולכסון של מטריצה}

תהי \L{%
\[
A\in\text{Mat}\left(3\times3,\mathbb{C}\right)
\]
}%
\L{%
\[
A=\begin{pmatrix}-1 & -2 & 2\\
-3 & -1 & 3\\
-3 & -2 & 4
\end{pmatrix}
\]
}%

)א( מצאו את כל הערכים העצמיים של \L{$A$} )ערכים עצמיים של מטריצה
פירושם ערכים עצמיים של ההעתקה הלינארית על \L{$\mathbb{C}^{3}$} שהמטריצה
מגדירה. זהו תרגיל סטנדרטי שיופיע בכל מועד של המבחן, ולכן חשוב לעשות
לפחות דוגמה אחת לא טריוויאלית(.

)ב( מצאו את הוקטורים העצמיים השייכים לערכים העצמיים.

)ג( מצאו מטריצה הפיכה \L{$Q$} כך שמתקיים \L{$A=QDQ^{-1}$} כאשר \L{$D$}
היא מטריצה אלכסונית.

\subsection{פתרון )א(}

ערך עצמי של \L{$A$} הוא כל פתרון \L{$\lambda\in\mathbb{C}$} של המשוואה
\L{%
\[
\det\left(A-\lambda I\right)=\begin{vmatrix}-1-\lambda & -2 & 2\\
-3 & -1-\lambda & 3\\
-3 & -2 & 4-\lambda
\end{vmatrix}=0
\]
}%

אפשר לפתח כאן את הדטרמיננטה, אבל אז נקבל פולינום ממעלה שלישית ויהיה
קשה למצוא את שורשיו )נצטרך לנחש איך לפרק אותו למכפלה(.

אנחנו רק צריכים לבדוק מתי היא מתאפסת.

כזכור, פעולות שורה אלמנטריות יכולות לשנות את הדטרמיננטה )למשל את הסימן
שלה או כפל בקבוע( אבל הן לא משנות את ההתאפסות שלה.

כלומר, הדטרמיננטה לפני פעולות אלו מתאפסת אם\textquotedblright ם הדטרמיננטה
לאחר פעולות אלו מתאפסת.

נזכיר שהסימון של המטריצה \L{$B$} לאחר שהופכים את השורות שלה לעמודות
הוא \L{$B^{T}$}.

למדנו בהרצאה \L{$29$} של לינארית 1 שמתקיים \L{%
\[
\det\left(B^{T}\right)=\det\left(B\right)
\]
}%

ולכן אותו הדבר נכון גם לגבי פעולות עמודה אלמנטריות.

חשוב לבצע את פעולות העמודה או השורה על המטריצה \L{$A-\lambda I$}
ולא קודם לדרג את \L{$A$} למשל ואז להחסיר את \L{$\lambda I$}, כי
דירוג של \L{$A$} משנה את הערכים העצמיים.

\L{%
\[
\begin{vmatrix}-1-\lambda & -2 & 2\\
-3 & -1-\lambda & 3\\
-3 & -2 & 4-\lambda
\end{vmatrix}\xrightarrow{R_{3}\to R_{3}-R_{2}}\begin{vmatrix}-1-\lambda & -2 & 2\\
-3 & -1-\lambda & 3\\
0 & \lambda-1 & 1-\lambda
\end{vmatrix}
\]
}%
\L{%
\begin{align*}
 & \underset{(1)}{=}\left(\lambda-1\right)\begin{vmatrix}-1-\lambda & -2 & 2\\
-3 & -1-\lambda & 3\\
0 & 1 & -1
\end{vmatrix}\\
 & \xrightarrow{C_{3}\to C_{3}+C_{2}}\left(\lambda-1\right)\begin{vmatrix}-1-\lambda & -2 & 0\\
-3 & -1-\lambda & 2-\lambda\\
0 & 1 & 0
\end{vmatrix}\\
 & \underset{(2)}{=}\left(\lambda-1\right)\left(0\cdot[]-1\cdot\begin{vmatrix}-1-\lambda & 0\\
-3 & 2-\lambda
\end{vmatrix}+0\cdot[]\right)\\
 & =\left(\lambda-1\right)\left(-1\cdot\left(-1-\lambda\right)\cdot\left(2-\lambda\right)-\left(-1\right)\cdot0\cdot\left(-3\right)\right)\\
 & =\left(\lambda-1\right)\left(1+\lambda\right)\left(2-\lambda\right)
\end{align*}
}%

\L{$(1)$} דטרמיננטה היא פונקציה מולטי-לינארית.

\L{$(2)$} פיתוח לפי השורה השלישית. בתרגילים כאלה כדאי לפתח לפי שורה
או עמודה עם הרבה אפסים.

\L{$\Leftarrow$} ערכים עצמיים של \L{$A$}: \L{%
\[
\lambda_{1}=1,\quad\lambda_{2}=-1,\quad\lambda_{3}=2
\]
}%

\subsection{פתרון )ב(}

מחפשים את כל הוקטורים \L{$v$} כך ש-\L{$Av=\lambda v$} כלומר \L{$\left(A-\lambda I\right)v=0$}
לכן נציב בכל פעם את הערך העצמי שנרצה לחשב את הוקטורים העצמיים שלו
ב-\L{$A-\lambda I$}.

\L{$\underline{\lambda_{1}=1}$}: \L{%
\[
\left(A-\lambda_{1}I\right)v=0
\]
}%
\L{%
\[
\begin{pmatrix}-2 & -2 & 2\\
-3 & -2 & 3\\
-3 & -2 & 3
\end{pmatrix}\begin{pmatrix}x\\
y\\
z
\end{pmatrix}=\begin{pmatrix}0\\
0\\
0
\end{pmatrix}
\]
}%

כזכור מלינארית 1, פש\textquotedblright א משמרות את מרחב הפתרונות.

שורות 2 ו-3 זהות אז שורה 3 לא נותנת מידע ואפשר לאפס אותה. \L{%
\[
\begin{pmatrix}-2 & -2 & 2\\
-3 & -2 & 3\\
0 & 0 & 0
\end{pmatrix}\begin{pmatrix}x\\
y\\
z
\end{pmatrix}=\begin{pmatrix}0\\
0\\
0
\end{pmatrix}
\]
}%

נוריד משורה 1 את שורה 2. \L{%
\[
\begin{pmatrix}1 & 0 & -1\\
-3 & -2 & 3\\
0 & 0 & 0
\end{pmatrix}\begin{pmatrix}x\\
y\\
z
\end{pmatrix}=\begin{pmatrix}0\\
0\\
0
\end{pmatrix}
\]
}%
\L{%
\[
\begin{cases}
x-z=0\\
-3x-2y+3z=0
\end{cases}
\]
}%
\L{%
\[
x=z,y=0
\]
}%

לכן הוקטורים העצמיים של \L{$\lambda_{1}$} הם הווקטורים שאינם אפס
השייכים לתמ\textquotedblright ו \L{%
\[
\left\{ \begin{pmatrix}x\\
0\\
x
\end{pmatrix}\middle|x\in\mathbb{C}\right\} =\text{Span}\left(\begin{pmatrix}1\\
0\\
1
\end{pmatrix}\right)=U_{1}
\]
}%

\L{$\underline{\lambda_{2}=-1}$}: \L{%
\[
\begin{pmatrix}0 & -2 & 2\\
-3 & 0 & 3\\
-3 & -2 & 5
\end{pmatrix}\begin{pmatrix}x\\
y\\
z
\end{pmatrix}=\begin{pmatrix}0\\
0\\
0
\end{pmatrix}
\]
}%
משורה \L{$1$}: \L{%
\[
y=z
\]
}%
משורה \L{$2$}: \L{%
\[
-x+z=0\Longrightarrow x=z
\]
}%
ובדיקה בשורה \L{$3$}: \L{%
\[
-3x-2y+5z=0\quad\checkmark
\]
}%

לכן הוקטורים העצמיים של \L{$\lambda_{2}$} הם הווקטורים שאינם אפס
השייכים לתמ\textquotedblright ו \L{%
\[
\left\{ \begin{pmatrix}x\\
x\\
x
\end{pmatrix}\middle|x\in\mathbb{C}\right\} =\text{Span}\left(\begin{pmatrix}1\\
1\\
1
\end{pmatrix}\right)=U_{2}
\]
}%

בדיקה: \L{$A\begin{pmatrix}1\\
1\\
1
\end{pmatrix}$} נותן בכל שורה \L{$-1$}, בדיוק \L{$-1\cdot\begin{pmatrix}1\\
1\\
1
\end{pmatrix}$}.

\L{$\underline{\lambda_{3}=2}$}: \L{%
\[
\begin{pmatrix}-3 & -2 & 2\\
-3 & -3 & 3\\
-3 & -2 & 2
\end{pmatrix}\begin{pmatrix}x\\
y\\
z
\end{pmatrix}=\begin{pmatrix}0\\
0\\
0
\end{pmatrix}
\]
}%
\L{%
\[
\begin{pmatrix}-3 & -2 & 2\\
-3 & -3 & 3\\
0 & 0 & 0
\end{pmatrix}\begin{pmatrix}x\\
y\\
z
\end{pmatrix}=\begin{pmatrix}0\\
0\\
0
\end{pmatrix}
\]
}%
\L{%
\[
\begin{pmatrix}0 & 1 & -1\\
-3 & -3 & 3\\
0 & 0 & 0
\end{pmatrix}\begin{pmatrix}x\\
y\\
z
\end{pmatrix}=\begin{pmatrix}0\\
0\\
0
\end{pmatrix}
\]
}%
משורה \L{$1$}: \L{%
\[
y=z
\]
}%
משורה \L{$2$}, עם \L{$y=z$}: \L{%
\[
-3x-3y+3z=-3x=0\Longrightarrow x=0
\]
}%

לכן הוקטורים העצמיים של \L{$\lambda_{3}$} הם הווקטורים שאינם אפס
השייכים לתמ\textquotedblright ו \L{%
\[
\left\{ \begin{pmatrix}0\\
x\\
x
\end{pmatrix}\middle|x\in\mathbb{C}\right\} =\text{Span}\left(\begin{pmatrix}0\\
1\\
1
\end{pmatrix}\right)=U_{3}
\]
}%

\subsection{פתרון )ג(}

נשתמש בטענה מהרצאה 3 המציינת כי אם ל-\L{$T_{A}$} יש \L{$n$} ערכים
עצמיים שונים זה מזה, אז \L{$A$} דומה למטריצה האלכסונית של הערכים
העצמיים, ומטריצת המעבר היא זו שעמודותיה הן הוקטורים העצמיים.

כלומר, \L{$Q$} היא המטריצה שעמודותיה הן הוקטורים העצמיים \L{%
\[
\begin{pmatrix}1 & 1 & 0\\
0 & 1 & 1\\
1 & 1 & 1
\end{pmatrix}
\]
}%

ו-\L{$D$} היא מטריצה אלכסונית של הערכים העצמיים התואמים \L{%
\[
\begin{pmatrix}1 & 0 & 0\\
0 & -1 & 0\\
0 & 0 & 2
\end{pmatrix}
\]
}%

\section{הערה - איברי האלכסון של מטריצה משולשית עליונה הם ערכים עצמיים}

ראינו בהרצאה הקודמת את המשפט לפיו מעל \L{$\mathbb{C}$}, כל אופרטור
ניתן לייצוג ע\textquotedblright י מטריצה משולשית עליונה. נזכיר:

יהי \L{$V$} מ\textquotedblright ו מעל \L{$\mathbb{C}$} כך ש- \L{%
\[
\dim V=n\geq1,\quad T\in\mathcal{L}\left(V\right)
\]
}%
אזי קיים בסיס \L{$\mathcal{B}_{V}$} ב-\L{$V$} כך ש- \L{%
\[
M\left(T,\mathcal{B}_{V},\mathcal{B}_{V}\right)=\begin{pmatrix}a_{1,1} & * & \cdots & *\\
0 & a_{2,2} & \cdots & *\\
\vdots &  & \ddots & \vdots\\
0 & 0 &  & a_{n,n}
\end{pmatrix}
\]
}%

נוסיף עכשיו:

ו- \L{$a_{1,1},\ldots,a_{n,n}$} הם ערכים עצמיים של \L{$T$}.

הם לא בהכרח שונים זה מזה, ייתכן שחלקם שווים.

אבל כל אחד מהם הוא ערך עצמי של \L{$T$}. 
\begin{proof}
נוכיח בעזרת הפולינום האופייני: \L{%
\[
\lambda I-M\left(T,\mathcal{B}_{V},\mathcal{B}_{V}\right)=\begin{pmatrix}\lambda-a_{1,1} & * & \cdots & *\\
0 & \lambda-a_{2,2} & \cdots & *\\
\vdots &  & \ddots & \vdots\\
0 & 0 &  & \lambda-a_{n,n}
\end{pmatrix}
\]
}%

זו מטריצה משולשית עליונה. נפתח לפי העמודה הראשונה: האיבר היחיד בה
שאינו אפס הוא \L{$\lambda-a_{1,1}$} במקום \L{$\left(1,1\right)$},
עם סימן \L{$\left(-1\right)^{1+1}=+1$}, והמינור שלו הוא שוב מטריצה
משולשית עליונה: \L{%
\[
\det\left(\lambda I-M\left(T,\mathcal{B}_{V},\mathcal{B}_{V}\right)\right)=\left(\lambda-a_{1,1}\right)\begin{vmatrix}\lambda-a_{2,2} & * & \cdots & *\\
0 & \lambda-a_{3,3} & \cdots & *\\
\vdots &  & \ddots & \vdots\\
0 & 0 &  & \lambda-a_{n,n}
\end{vmatrix}
\]
}%
\L{%
\[
=\left(\lambda-a_{1,1}\right)\left(\lambda-a_{2,2}\right)\begin{vmatrix}\lambda-a_{3,3} & * & \cdots & *\\
0 & \lambda-a_{4,4} & \cdots & *\\
\vdots &  & \ddots & \vdots\\
0 & 0 &  & \lambda-a_{n,n}
\end{vmatrix}
\]
}%
\L{%
\[
=\dots=\left(\lambda-a_{1,1}\right)\left(\lambda-a_{2,2}\right)\cdots\left(\lambda-a_{n,n}\right)
\]
}%
)זו ההוכחה של הטענה מלינארית 1 שדטרמיננטה של מטריצה משולשית היא מכפלת
איברי האלכסון(

מכאן השורשים הם בדיוק \L{$a_{1,1},\ldots,a_{n,n}$}. 
\end{proof}
יש לנו עד עכשיו שתי תוצאות מרכזיות:

)1( מעל כל שדה, אם ל-\L{$T$} יש \L{$n$} ערכים עצמיים שונים, אז \L{$T$}
לכסינה.

)2( מעל \L{$\mathbb{C}$}, לכל \L{$T$} קיים בסיס שבו המטריצה משולשית
עליונה.

השאלה: מה אפשר לומר מעל שדה כללי, ומתי מגיעים למטריצה אלכסונית?

לשם כך צריך ללמוד פולינומים שקשורים להעתקה, ובפרט את הפולינום המינימלי.

\part{הפולינום המינימלי}

\section{נוטציה: \L{$p(T)$}}

יהי \L{$V$} מ\textquotedblright ו ויהי \L{$T\in\mathcal{L}(V)$}.

בנוסף, יהי \L{$p\in\mathcal{P}(\mathbb{F})$} הפולינום \L{%
\[
\forall z\in\mathbb{F}:p(z)=a_{0}+a_{1}z+a_{2}z^{2}+\dots+a_{m}z^{m}
\]
}%

אז \L{$p(T)$} הוא האופרטור על \L{$V$} המוגדר ע\textquotedblright י
\L{%
\[
p(T)=a_{0}I+a_{1}T+a_{2}T^{2}+\dots+a_{m}T^{m}
\]
}%

כאשר \L{$T^{k}=\underset{k\text{ times}}{\underbrace{T\circ\dots\circ T}}$}
)אפשר גם להגדיר \L{$T^{0}=I_{V}$} וזה מקל על הכתיבה( ומתקיים \L{%
\[
p(T)(v)=a_{0}\cdot I(v)+a_{1}\cdot T(v)+a_{2}\cdot T\left(T\left(v\right)\right)+\dots
\]
}%

הרכבה, חיבור וכפל בסקלר של ה\textquotedblright ל נותנים ה\textquotedblright ל
לכן גם התוצאה היא אופרטור לינארי על \L{$V$} וזה מוגדר היטב.

\section{פולינום מתוקן}

פולינום מתוקן הוא פולינום בו המקדם המוביל הוא 1.

\section{הגדרה - פולינום מינימלי של \L{$T$}}

יהי \L{$V$} מ\textquotedblright ו נ\textquotedblright ס מעל שדה \L{$\mathbb{F}$}
ויהי \L{$T\in\mathcal{L}(V)$}.

אז הפולינום המינימלי של \L{$T$} הוא הפולינום המתוקן \L{$p\in\mathcal{P}(\mathbb{F})$}
מהדרגה הנמוכה ביותר כך ש-\L{$p(T)=0$} )אופרטור האפס(: \L{%
\[
\forall v\in V:p(T)(v)=a_{0}v+a_{1}T(v)+\dots+\underset{=1}{a_{m}}T^{m}(v)=0
\]
}%

מהדרגה הנמוכה ביותר, דהיינו: \L{%
\[
\forall\ 0\ne q\in\mathcal{P}(\mathbb{F}):q(T)=0\Rightarrow\deg p\le\deg q
\]
}%

בהרצאה יבגני גם דרש \L{$\dim(V)=n\ge1$} בשביל לא להתעסק עם המקרה
הטריוויאלי.

כתבנו \textbf{ה}פולינום \textbf{ה}מינימלי ולכן אחרי ההגדרה נצטרך להוכיח
קיום ויחידות.

על השאלה הרלוונטית, שהיא איך זה קשור לערכים עצמיים, נענה בהרצאה הבאה.

\section{משפט - קיום ויחידות של הפולינום המינימלי}

יהי \L{$V$} מ\textquotedblright ו נ\textquotedblright ס מעל \L{$\mathbb{F}$}
עם \L{$\dim(V)=n\ge1$} ויהי \L{$T\in\mathcal{L}(V)$} )אקסלר לא מניח
שהמימד גדול שווה לאחד ומטפל בנפרד במקרה הטריוויאלי בו \L{$I$} הוא
אופרטור האפס ו-\L{$p$} הוא הפולינום הקבוע 1(.

אז קיים פולינום מתוקן \L{$p\in\mathcal{P}(\mathbb{F})$} כך ש-\L{$p(T)=0$}
וגם \L{%
\[
\forall\ 0\ne q\in\mathcal{P}(\mathbb{F}):q(T)=0\Rightarrow\deg p\le\deg q
\]
}%

והוא יחיד.

\subsection{למה - \L{$\mathcal{L}\left(V\right)$} הוא מ\textquotedblright ו
עם \L{$\dim\left(\mathcal{L}\left(V\right)\right)=n^{2}$}}

)יבגני לא הוכיח זאת אלא התייחס לזה בתור טענה מלינארית 1. לא מצאתי
אותה ולכן אני מוכיח.(

יהי \L{$V$} מ\textquotedblright ו נ\textquotedblright ס מעל \L{$\mathbb{F}$}
עם \L{$\dim(V)=n$}. אז לכל \L{$S,T\in\mathcal{L}(V)$}, \L{$v\in V$}
ו-\L{$c\in\mathbb{F}$} נגדיר \L{%
\[
\left(T+S\right)(v)=T(v)+S(v),\quad\left(cT\right)(v)=c\cdot T(v)
\]
}%

ניתן לראות )לא אוכיח זאת כאן( ש-\L{$T+S,cT\in\mathcal{L}(V)$} ושהקבוצה
\L{$\mathcal{L}(V)$} יחד עם הגדרות אלה של חיבור וכפל בסקלר מקיימת
את האקסיומות של מ\textquotedblright ו.

\L{$V$} מ\textquotedblright ו נ\textquotedblright ס ולכן יש לו בסיס
\L{$\mathcal{B}_{V}=\left(v_{1},\dots,v_{n}\right)$}.

נגדיר \L{$\Phi:\mathcal{L}(V)\rightarrow\text{Mat}\left(n\times n,\mathbb{F}\right)$}
על ידי \L{$\Phi(T)=M\left(T,\mathcal{B}_{V},\mathcal{B}_{V}\right)$}.

אראה סקיצה של ההוכחה כיוון שזה שייך ללינארית 1.

\L{$\Phi$} לינארית: \L{%
\[
\forall1\le j\le n:\left[\left(T+S\right)\left(v_{j}\right)\right]_{\mathcal{B}_{V}}=\left[T\left(v_{j}\right)+S\left(v_{j}\right)\right]_{\mathcal{B}_{V}}\underset{(1)}{=}\left[T\left(v_{j}\right)\right]_{\mathcal{B}_{V}}+\left[S\left(v_{j}\right)\right]_{\mathcal{B}_{V}}
\]
}%
\L{%
\[
\forall1\le j\le n:\left[\left(cT\right)\left(v_{j}\right)\right]_{\mathcal{B}_{V}}=\left[c\cdot T\left(v_{j}\right)\right]_{\mathcal{B}_{V}}\underset{(1)}{=}c\cdot\left[T\left(v_{j}\right)\right]_{\mathcal{B}_{V}}
\]
}%
\L{%
\[
\therefore\Phi(S+T)=\Phi(S)+\Phi(T),\quad\Phi(cT)=c\Phi(T)
\]
}%

\L{$(1)$} למדנו בהרצאה 11 של לינארית 1 שוקטור הקואורדינטות משמר חיבור
וכפל בסקלר.

\L{$\Phi$} חח\textquotedblright ע: למדנו בלינארית 1 שה\textquotedblright ל
נקבעת ע\textquotedblright י המטריצה המייצגת שלה ולכן \L{$\Phi(T)=\Phi(S)\Rightarrow T=S$}.

\L{$\Phi$} על: למדנו בלינארית 1 שכל מטריצה ניתנת לייצוג ע\textquotedblright י
ה\textquotedblright ל ולכן לכל \L{$A\in\text{Mat}\left(n\times n,\mathbb{F}\right)$}
קיים \L{$T\in\mathcal{L}(V)$} כך ש-\L{$\Phi(T)=A$}.

מכאן \L{$\Phi$} היא איזומורפיזם.

הבסיס הסטנדרטי של \L{$\text{Mat}\left(n\times n,\mathbb{F}\right)$}
הוא \L{$E_{1,1},\dots,E_{n,n}$} לכן \L{$\dim\left(\text{Mat}\left(n\times n,\mathbb{F}\right)\right)=n^{2}$}.

אז לפי משפט שלמדנו בלינארית 1 היות ו-\L{$\Phi$} איזומורפיזם, לתחום
ולטווח שלה אותו מימד דהיינו \L{$\dim\left(\mathcal{L}\left(V\right)\right)=n^{2}$}.

\subsection{למה - \L{$\Phi:\mathcal{L}(V)\rightarrow\text{Mat}\left(n\times n,\mathbb{F}\right)$}
מכבדת כפל}

)גם את זה יבגני לא הוכיח(

לפי הרצאה \L{$19$} בלינארית 1, המטריצה של הרכבה היא מכפלת המטריצות.
מכאן \L{%
\[
\Phi\left(ST\right)=\Phi(S)\Phi(T)
\]
}%

\subsection{הוכחת המשפט}
\begin{proof}
\textbf{1. הוכחת קיום:}

לפי הלמה, \L{$\mathcal{L}\left(V\right)$} הוא מ\textquotedblright ו
עם \L{$\dim\left(\mathcal{L}\left(V\right)\right)=n^{2}$}. נסתכל
על הרשימה \L{%
\[
\left(I,T,T^{2},\ldots,T^{n^{2}}\right)
\]
}%
ברשימה זאת יש \L{$n^{2}+1$} איברים במרחב ממימד \L{$n^{2}$}, ולכן
לפי משפט מלינארית 1 היא תלויה לינארית במרחב \L{$\mathcal{L}\left(V\right)$}.

לפי למת התלות הלינארית \L{%
\[
\exists1\underset{(*)}{\le}m\le n^{2}:T^{m}\in\text{Span}\left(I,T,\ldots,T^{m-1}\right)
\]
}%

\L{$(*)$} הלמה מבטיחה במקרה זה רק \L{$m\ge0$}, אבל \L{$n\ge1$}
לכן \L{$I\ne0$} ועם \L{$m=0$} מקבלים סתירה.

נבחר את ה-\L{$m$} הקטן ביותר המקיים זאת. אז הרשימה \L{$\left(I,T,\ldots,T^{m-1}\right)$}
בת\textquotedblright ל: אילו הייתה תלויה, מכיוון ש-\L{$I\neq0$} למת
התלות הלינארית הייתה נותנת \L{$k<m$} עם \L{$T^{k}\in\text{Span}\left(I,\ldots,T^{k-1}\right)$},
בסתירה למינימליות של \L{$m$}.

אז קיימים \L{$a_{0},\dots,a_{m-1}\in\mathbb{F}$} כך ש- \L{%
\[
a_{0}I+a_{1}T+\dots+a_{m-1}T^{m-1}=T^{m}
\]
}%
\L{%
\[
\therefore a_{0}I+a_{1}T+\dots+a_{m-1}T^{m-1}+T^{m}=0
\]
}%
ונגדיר את הפולינום הבא: \L{%
\[
p(z)=a_{0}+a_{1}z+\dots+a_{m-1}z^{m-1}+z^{m}
\]
}%
אז לכל \L{$v\in V$} \L{%
\[
p(T)(v)=a_{0}I(v)+a_{1}T(v)+\dots+a_{m-1}T^{m-1}(v)+T^{m}(v)=0(v)=0
\]
}%

בנוסף, היות והרשימה בת\textquotedblright ל, עבור \L{$b_{0},\dots,b_{m-1}\in\mathbb{F}$}
כך ש- \L{%
\[
b_{0}I+b_{1}T+\dots+b_{m-1}T^{m-1}=0
\]
}%
מתקיים \L{%
\[
b_{0}=\dots=b_{m-1}=0
\]
}%
כלומר, אם \L{$q$} הוא פולינום שונה מאפס שעבורו מתקיים \L{$q(T)=0$},
הדרגה שלו חייבת להיות גדולה מ-\L{$m-1$} )אחרת לפי מה שהראינו עכשיו
הוא פולינום האפס וזאת סתירה( ולכן גדולה שווה מ-\L{$m$}, כלומר \L{$\deg(q)\ge\deg(p)$}
כנדרש.

\textbf{2. הוכחת יחידות:}

נניח כי קיים פולינום מינימלי נוסף \L{%
\[
\widetilde{p}(z)=b_{0}+b_{1}z+\dots+b_{m-1}z^{m-1}+z^{m}
\]
}%
ראשית, יש להם אותה מעלה, כי מהגדרת הפולינום המינימלי \L{$\deg(p)\ge\deg(\widetilde{p})$}
וגם \L{$\deg(\widetilde{p})\ge\deg(p)$}.

שנית, שניהם מתוקנים, ולכן המקדם של \L{$z^{m}$} אצל שניהם הוא 1 אז
כשמחסרים אחד מהשני הוא נעלם. \L{%
\[
\left(a_{0}-b_{0}\right)I+\left(a_{1}-b_{1}\right)T+\dots+\left(a_{m-1}-b_{m-1}\right)T^{m-1}=\left(p-\widetilde{p}\right)\left(T\right)=0
\]
}%
אבל כפי שהוכחנו, \L{$\left(I,T,\ldots,T^{m-1}\right)$} היא רשימה
בת\textquotedblright ל, אז \L{%
\begin{align*}
a_{0}-b_{0} & =0\\
 & \vdots\\
a_{m-1}-b_{m-1} & =0
\end{align*}
}%
\L{%
\[
\therefore a_{0}=b_{0},\dots,a_{m-1}=b_{m-1}
\]
}%
כלומר \L{$\widetilde{p}=p$} כנדרש. 
\end{proof}
מבחינה מעשית, אין שיטה חכמה כללית למצוא נוסחה מפורשת לפולינום המינימלי;
גם עבור \L{$3\times3$} זה יכול להיות קשה.

השיטה שראינו עובדת לכל העתקה במ\textquotedblright ו נ\textquotedblright ס:
בוחרים בסיס, עוברים למטריצה, ומחפשים את המעלה הקטנה ביותר שבה קיים
פולינום מתוקן שמאפס את המטריצה. אם ההעתקה פשוטה, זה קל.

\section{דוגמה - הפולינום המינימלי של \L{$T\left(z_{1},z_{2}\right)=\left(-z_{2},z_{1}\right)$}}

תהי \L{$T\in\mathcal{L}\left(\mathbb{C}^{2}\right)$} הבאה \L{%
\[
T\left(\left(z_{1},z_{2}\right)\right)=\left(-z_{2},z_{1}\right)
\]
}%
מהו הפולינום המינימלי של \L{$T$}?

\subsection{פתרון}

ניקח את הבסיס הסטנדרטי \L{%
\[
\mathcal{B}_{V}=\left(\left(1,0\right),\left(0,1\right)\right)=\left(v_{1},v_{2}\right)
\]
}%

נפעיל עליו את האופרטור \L{%
\[
T\left(\left(1,0\right)\right)=\left(0,1\right)=0\cdot\left(1,0\right)+1\cdot\left(0,1\right)
\]
}%
\L{%
\[
T\left(\left(0,1\right)\right)=\left(-1,0\right)=-1\cdot\left(1,0\right)+0\cdot\left(0,1\right)
\]
}%
\L{%
\[
M\left(T,\mathcal{B}_{V},\mathcal{B}_{V}\right)=\begin{pmatrix}0 & -1\\
1 & 0
\end{pmatrix}
\]
}%

דיברנו כבר על ההעתקה הזאת.

לפי המשפט, מספיק למצוא פולינום מתוקן ממעלה מינימלית שמאפס את המטריצה;
הוא יהיה הפולינום המינימלי, לפי היחידות.

מתחילים ממעלה \L{$1$} ועולים.

\L{$\underline{m=1}$}: נחפש \L{$a_{0}\in\mathbb{C}$} כך ש- \L{%
\[
a_{0}\begin{pmatrix}1 & 0\\
0 & 1
\end{pmatrix}+1\cdot\begin{pmatrix}0 & -1\\
1 & 0
\end{pmatrix}=\begin{pmatrix}0 & 0\\
0 & 0
\end{pmatrix}
\]
}%
\L{%
\[
\begin{pmatrix}a_{0} & -1\\
1 & a_{0}
\end{pmatrix}=\begin{pmatrix}0 & 0\\
0 & 0
\end{pmatrix}
\]
}%
ובגלל ה-\L{$-1$} וה-\L{$1$} זה לא ייתכן. לכן אין פולינום ממעלה \L{$1$}
שמאפס את \L{$T$}, והפולינום המינימלי הוא ממעלה גבוהה יותר.

\L{$\underline{m=2}$}: נחפש \L{$a_{0},a_{1}\in\mathbb{C}$} כך ש-
\L{%
\[
a_{0}\begin{pmatrix}1 & 0\\
0 & 1
\end{pmatrix}+a_{1}\begin{pmatrix}0 & -1\\
1 & 0
\end{pmatrix}+1\cdot\begin{pmatrix}0 & -1\\
1 & 0
\end{pmatrix}^{2}=\begin{pmatrix}0 & 0\\
0 & 0
\end{pmatrix}
\]
}%
נשים לב שמתקיים \L{%
\[
\begin{pmatrix}0 & -1\\
1 & 0
\end{pmatrix}^{2}=\begin{pmatrix}0 & -1\\
1 & 0
\end{pmatrix}\begin{pmatrix}0 & -1\\
1 & 0
\end{pmatrix}=-I
\]
}%
אז למעשה אנו מחפשים \L{$a_{0},a_{1}\in\mathbb{C}$} כך ש- \L{%
\[
a_{0}\begin{pmatrix}1 & 0\\
0 & 1
\end{pmatrix}+a_{1}\begin{pmatrix}0 & -1\\
1 & 0
\end{pmatrix}+1\cdot\left(-1\right)\begin{pmatrix}1 & 0\\
0 & 1
\end{pmatrix}=\begin{pmatrix}0 & 0\\
0 & 0
\end{pmatrix}
\]
}%
\L{%
\[
\left(a_{0}-1\right)\begin{pmatrix}1 & 0\\
0 & 1
\end{pmatrix}+a_{1}\begin{pmatrix}0 & -1\\
1 & 0
\end{pmatrix}=\begin{pmatrix}0 & 0\\
0 & 0
\end{pmatrix}
\]
}%
\L{%
\[
\begin{pmatrix}a_{0}-1 & 0\\
0 & a_{0}-1
\end{pmatrix}+\begin{pmatrix}0 & -a_{1}\\
a_{1} & 0
\end{pmatrix}=\begin{pmatrix}0 & 0\\
0 & 0
\end{pmatrix}
\]
}%
\L{%
\[
\begin{pmatrix}a_{0}-1 & -a_{1}\\
a_{1} & a_{0}-1
\end{pmatrix}=\begin{pmatrix}0 & 0\\
0 & 0
\end{pmatrix}
\]
}%

מקבלים \L{$a_{0}=1,a_{1}=0$}.

מכאן, הפולינום המינימלי של \L{$T$} הוא \L{%
\[
p(z)=1+z^{2}
\]
}%

מצאנו פולינום מתוקן ממעלה \L{$2$} שמתאפס על המטריצה.

השלב הבא הוא להבין שאם זה מתקיים עבור המטריצה, זה מתקיים גם עבור ההעתקה,
כי המטריצה מייצגת את ההעתקה.

אם לא היינו מצליחים במעלה \L{$2$}, היינו ממשיכים למעלה \L{$3$} וכן
הלאה.

\subsection{הערה - שורשי הפולינום המינימלי וערכים עצמיים}

שימו לב לשורשים של \L{$z^{2}+1$}: הם \L{$i$} ו-\L{$-i$}.

אלה בדיוק הערכים העצמיים של \L{$T$} שמצאנו בעבר.

בדוגמה הזאת יש קשר בין הפולינום המינימלי לערכים העצמיים, וההשערה היא
שזה תמיד כך: כל שורש של הפולינום המינימלי של \L{$T$} הוא ערך עצמי
של \L{$T$}, ולהפך, כל ערך עצמי של \L{$T$} הוא שורש של הפולינום המינימלי.

זה באמת נכון, ונוכיח זאת בהמשך.

בנוסף, גם \L{$\det\left(\lambda I-M\right)$} הוא פולינום )הפולינום
האופייני(, ומתברר שאם מציבים בו את \L{$T$} במקום \L{$\lambda$} מקבלים
\L{$0$}.

זה משפט \L{$\text{Cayley-Hamilton}$} )הידוע לשמצה(, וגם אותו נוכיח
בהמשך.

\section{למה - הפולינום המינימלי של אופרטור שווה לפולינום המינימלי של המטריצה
המייצגת שלו )לא היה בהרצאה(}

יהי \L{$V$} מ\textquotedblright ו מעל \L{$\mathbb{F}$} עם \L{$\dim V=n\geq1$}
ויהי \L{$T\in\mathcal{L}\left(V\right)$}.

יהי \L{$\mathcal{B}_{V}$} בסיס של \L{$V$}. נסמן \L{%
\[
M=M\left(T,\mathcal{B}_{V},\mathcal{B}_{V}\right)
\]
}%
אז הפולינום המינימלי של \L{$T$} שווה לפולינום המינימלי של \L{$M$}. 
\begin{proof}
1. נסמן \L{$\Phi:\mathcal{L}\left(V\right)\to\text{Mat}\left(n\times n,\mathbb{F}\right)$},
\L{$\Phi\left(S\right)=M\left(S,\mathcal{B}_{V},\mathcal{B}_{V}\right)$}.

לפי הלמות, \L{$\Phi$} היא איזומורפיזם ומכבדת כפל.

בנוסף \L{$\Phi\left(I\right)=I_{n}$}.

2. נוכיח שלכל פולינום \L{$q\left(z\right)=a_{0}+a_{1}z+\dots+a_{k}z^{k}$}
מתקיים \L{%
\[
\Phi\left(q\left(T\right)\right)=q\left(M\right)
\]
}%
מהכפליות, באינדוקציה על \L{$j$}, \L{%
\[
\forall j\ge0:\Phi\left(T^{j}\right)=\Phi\left(T\right)^{j}=M^{j}
\]
}%
לכן \L{%
\begin{align*}
\Phi\left(q\left(T\right)\right) & =\Phi\left(a_{0}I+a_{1}T+\dots+a_{k}T^{k}\right)\\
 & =a_{0}\Phi\left(I\right)+a_{1}\Phi\left(T\right)+\dots+a_{k}\Phi\left(T^{k}\right)\\
 & =a_{0}I_{n}+a_{1}M+\dots+a_{k}M^{k}\\
 & =q\left(M\right)
\end{align*}
}%

3. נוכיח שלכל פולינום \L{$q$} מתקיים \L{%
\[
q\left(T\right)=0\iff q\left(M\right)=0
\]
}%
\L{$\Phi$} ה\textquotedblright ל ולכן \L{$\Phi\left(0\right)=0$}.

\L{$\Phi$} חח\textquotedblright ע ולכן הגרעין שלה הוא אפס כלומר \L{$\Phi\left(X\right)=0$}
אם\textquotedblright ם \L{$X=0$}.

מכאן \L{%
\[
q\left(T\right)=0\iff\Phi\left(q\left(T\right)\right)=0\underset{(2)}{\iff}q\left(M\right)=0
\]
}%

4. יהי \L{$p$} הפולינום המינימלי של \L{$T$}. נראה ש-\L{$p$} מקיים
את שני תנאי ההגדרה עבור \L{$M$}.

\L{$p$} מתוקן, ו-\L{$p\left(T\right)=0$}, ולכן לפי שלב 3 גם \L{$p\left(M\right)=0$}.

יהי \L{$0\neq q$} פולינום עם \L{$q\left(M\right)=0$}.

לפי שלב 3, \L{$q\left(T\right)=0$}.

\L{$p$} מינימלי עבור \L{$T$}, ולכן \L{$\deg p\leq\deg q$}.

לכן \L{$p$} הוא פולינום מינימלי של \L{$M$}, ומיחידות הפולינום המינימלי
הוא הפולינום המינימלי של \L{$M$}.

5. הכיוון ההפוך זהה עם החלפת התפקידים של \L{$T$} ו-\L{$M$}, כי שלב
3 הוא שקילות. 
\end{proof}

\end{document}
